\section{Discussion}\label{sec:disc}

\noindent\textbf{What Coordination Changes in Practice.} The empirical evidence demonstrates that introducing stateful, bounded-authority agents genuinely alters the behavioral dynamics of the interview preparation system. By tracking conceptual weaknesses across turns, the Strategy Agent smooths difficulty transitions, reducing high-frequency level oscillations (estimated reduction 2.67, with the cluster-bootstrap 95\% interval of $[-4.50, -0.24]$ entirely below zero) and directional reversals (reduction 1.88, cluster-bootstrap 95\% CI $[-3.12, -0.27]$) compared to the frozen single-policy PPO baseline. However, scientific rigor demands that we distinguish behavioral stabilization from pedagogical quality. Replay on fixed interview responses cannot determine whether a smoother difficulty progression promotes superior candidate learning, retention, or interview confidence. A smoother trajectory demonstrates algorithmic consistency, not verified educational superiority.

\noindent\textbf{The Limits and Realities of Agent Specialization.} A pervasive narrative in the multi-agent literature assumes that decomposing complex reasoning into specialized agents inherently produces superior task performance~\cite{autogen,metagpt}. Our findings challenge this assumption in deterministic settings. Under the tested configuration, specialization produces \textbf{zero algorithmic divergence}: decision outputs and difficulty trajectories are identical to an ablated single-agent or unified configuration ($M_3 \equiv M_1$, 160/160 identical sessions, $\mathbb{P}(\text{agree}) = 1.0$ across 1{,}280 turns). 

Instead, specialization functions primarily as an architectural decomposition with specific design characteristics rather than measured empirical productivity outcomes:
\begin{enumerate}
\item \emph{Lifecycle Decomposition (Design Characteristic):} Strategy and Coaching address distinct concerns (macro-level progression vs micro-level pedagogical presentation). While we do not measure developer productivity empirically, role separation provides explicit boundary contracts and independent schemas.
\item \emph{Decoupled State Ledgers (Architectural Property):} Distributing state across dedicated private ledgers maintains isolated schema boundaries for difficulty history and concept tracking, enabling independent auditability without coupling domain models.
\item \emph{Operational Overhead (Measured Empirical Cost):} Modularity incurs measured runtime costs across distinct scopes: $+1.04$\,ms median end-to-end orchestration overhead per turn (encompassing agent invocation, peer messaging, and arbitration; 47.7 calls per session), alongside narrower serialization micro-benchmark differences ($0.0092$\,ms vs $0.0038$\,ms) and component dispatch latency ($0.0440$\,ms vs $0.0401$\,ms).
\end{enumerate}

\noindent\textbf{Bounded Authority vs Emergent Multi-Agent Intelligence.} Our architecture intentionally eschews open-ended agent negotiation or decentralized consensus mechanisms. In high-stakes educational and evaluation settings, unconstrained agent communication introduces risks of specification drift, circular reasoning, and cascading hallucination~\cite{mast}. By enforcing a centralized orchestrator as the sole system of record, restricting agents to typed proposal contracts, and applying strict threshold-gated arbitration, our system realizes the principle of least privilege~\cite{ss75}. The system achieves predictable, verifiable coordination without sacrificing centralized safety bounds.

\noindent\textbf{Relevance to ADCoM 2027.} This investigation aligns directly with the core themes of the 31st International Conference on Advanced Computing and Communications (ADCoM 2027), particularly Track 2: \emph{Agent-Native Development and Lifecycle Engineering}. By rigorously auditing contract boundaries, quantifying the trade-offs of agent specialization, and isolating coordination mechanisms through cluster-bootstrap replay, this work demonstrates how multi-agent architectures can be engineered, bounded, and evaluated with empirical integrity.
